\chapter{CONCLUSÃO}

A consolidação das redes de acesso óptico baseadas na arquitetura \gls{pon} representa um pilar fundamental para o atendimento à crescente demanda por conectividade de alta velocidade em redes \gls{ftth}. Contudo, a complexidade inerente ao dimensionamento físico da camada passiva — que exige a contabilização rigorosa da atenuação da fibra óptica, conectores, emendas por fusão e elementos de divisão óptica (\textit{splitters}) — impõe severos desafios aos projetistas e engenheiros de telecomunicações. Este capítulo sintetiza as principais realizações do trabalho, destacando o cumprimento dos objetivos, as contribuições teóricas e práticas, as limitações identificadas e as perspectivas de trabalhos futuros.

\section{Síntese do Trabalho}

O problema que motivou a realização desta pesquisa reside na histórica fragmentação dos processos de planejamento e dimensionamento de redes \gls{pon}. Tradicionalmente, o projeto de infraestrutura óptica é conduzido através de ferramentas desacopladas: a representação visual e cartográfica da topologia é realizada em \textit{softwares} de desenho esquemático ou ambiente \gls{gis}, enquanto o cálculo do orçamento óptico é mantido em planilhas eletrônicas independentes ou notas de cálculo manuais. Essa desconexão impõe alto volume de retrabalho manual e introduz grande suscetibilidade a falhas humanas, tais como o descompasso entre a quantidade real de fusões do enlace e os parâmetros adotados na planilha, culminando no dimensionamento incorreto e em clientes com nível de potência inoperante.

Diante dessa problemática, estabeleceu-se a seguinte pergunta central de pesquisa: \textit{como integrar a modelagem visual da topologia ao cálculo do orçamento óptico de modo que as alterações do projeto produzam indicadores técnicos em tempo real para apoiar a tomada de decisão do projetista?}

Para responder a essa questão, percorreu-se um itinerário metodológico que compreendeu:
\begin{enumerate}
    \item \textbf{Levantamento de Requisitos e Análise Correlata:} Investigação abrangente das normas técnicas internacionais (série ITU-T G.984), fundamentos do orçamento óptico e mapeamento das lacunas existentes em \textit{softwares} de mercado e da literatura especializada (Capítulos 2 e 3);
    \item \textbf{Engenharia de Requisitos e Arquitetura do Sistema:} Modelação de um motor de cálculo óptico em Python baseado na teoria dos grafos e travessia em árvore, com especificação de uma \gls{api} \gls{rest} leve construída com FastAPI (Capítulos 4 e 5);
    \item \textbf{Desenvolvimento da Interface Interativa:} Construção do \textit{frontend} web utilizando React e D3.js, proporcionando manipulação vetorial da topologia via \textit{drag-and-drop}, atribuição visual de parâmetros e atualização bidirecional de sinal em tempo real;
    \item \textbf{Validação Empírica e Estudo de Casos:} Realização de testes funcionais e de desempenho sob dois cenários práticos reais — o Teste 1 em ambiente urbano denso (Bairro Cidade Alta em Videira/SC) e o Teste 2 em ambiente rural com distribuição desbalanceada de longo alcance (Município de Ibirama/SC, composto por 14 Caixas de Terminação Óptica — \glspl{cto}).
\end{enumerate}

Os resultados obtidos demonstraram que a aplicação web desenvolvida, denominada \textbf{PON Planner}, resolve com precisão a fragilidade do cálculo manual. O motor de cálculo respondeu em menos de 10 ms para topologias com até 128 \glspl{onu}, oferecendo retracagem instantânea do sinal óptico a cada modificação nos nós ou enlaces. Além disso, o teste comparativo em Ibirama/SC revelou que o cálculo automatizado identifica atenuações acumuladas (de $+0,7$ dB a $+1,3$ dB) negligenciadas nas contas manuais, permitindo diagnosticar antecipadamente \glspl{cto} inoperantes (sinal inferior a $-27$ dBm) e reduzindo o tempo de balanceamento da rede de mais de 1,5 hora para apenas 5 minutos.

No que tange à validação dos objetivos, confirma-se formalmente que o **Objetivo Geral** e os **quatro Objetivos Específicos** foram integralmente alcançados:

\begin{itemize}
    \item \textbf{Objetivo Específico 1 (Modelagem Visual):} Alcançado mediante a implementação do \textit{canvas} interativo com React e D3.js, permitindo a inserção, edição, exclusão e movimentação de \glspl{olt}, \glspl{ceo}, \glspl{cto}, \textit{splitters} e \glspl{onu}.
    \item \textbf{Objetivo Específico 2 (Cálculo do Orçamento Óptico):} Alcançado através do módulo \texttt{calculo\_optico.py}, que realiza a navegação recursiva no grafo direcionado da rede e calcula a atenuação total em 1490 nm, considerando perdas de fibra por quilômetro, conectores nas extremidades, emendas por fusão e razões de divisão ópticas (balanceadas de $1:2$ a $1:16$ e desbalanceadas de $10/90$ a $50/50$).
    \item \textbf{Objetivo Específico 3 (Indicadores Técnicos e \textit{Status} Visual):} Alcançado pelo mecanismo de atualização em tempo real de potência e margem óptica, sinalizado visualmente através do código de cores dos nós (verde para níveis \textit{OK} entre $-8$ dBm e $-25$ dBm, amarelo para \textit{Atenção} entre $-25$ dBm e $-27$ dBm, e vermelho para \textit{Inoperante} abaixo de $-27$ dBm ou acima de $-8$ dBm).
    \item \textbf{Objetivo Específico 4 (Persistência e Interoperabilidade \gls{json}):} Alcançado mediante a especificação da estrutura de dados em formato \gls{json}, permitindo salvar a estrutura completa da topologia em arquivo local, recuperar projetos previamente armazenados e exportar relatórios de validação sem perda de integridade.
\end{itemize}

\section{Contribuições}

O trabalho oferece contribuições significativas para o campo de planejamento de redes e engenharia de telecomunicações, as quais podem ser categorizadas sob quatro perspectivas:

\begin{enumerate}
    \item \textbf{Contribuições Teóricas:} Proposição e formalização de um modelo gráfico representativo de arquiteturas \gls{pon} como grafos direcionados acíclicos ponderados. O modelo integra a física de atenuação óptica (comprimento de onda, atenuação específica da fibra G.652.D, perda de inserção de acopladores e conectores) à teoria dos grafos, permitindo calcular analiticamente o orçamento óptico \textit{downstream} de qualquer topologia ponto-multiponto de forma matricial e recursiva.
    
    \item \textbf{Contribuições Práticas:} Desenvolvimento e disponibilização da aplicação \textbf{PON Planner} em código aberto (\textit{open-source}), devidamente containerizada através de Docker. A ferramenta elimina a dependência de licenças computacionais comerciais de alto custo para pequenos e médios Provedores de Serviços de Internet (ISPs) e oferece um ambiente funcional acessível para equipes de engenharia, projetistas e consultorias técnicas.
    
    \item \textbf{Contribuições Metodológicas:} Estabelecimento de um fluxo unificado de projeto (\textit{single source of truth}), onde a alteração dos componentes na interface gráfica reflete instantaneamente nos parâmetros de transmissão física de toda a árvore \textit{downstream}. Essa abordagem substitui a metodologia tradicional desacoplada (desenho em \textit{software} vetorial e cálculo em planilha), garantindo rastreabilidade e reduzindo dramaticamente os erros de digitação e dimensionamento.
    
    \item \textbf{Contribuições para a Área:} O avanço promovido por esta pesquisa pode ser sintetizado em três aspectos centrais: (1) integração em tempo real entre a modelagem visual de topologias \gls{ftth} e o cálculo automatizado de transmissão óptica; (2) disponibilização de uma ferramenta leve e aberta de apoio à decisão para engenheiros de telecomunicações, projetistas e provedores de serviços; e (3) consolidação de uma estratégia visual para o diagnóstico prévio de viabilidade técnica em topologias desbalanceadas rurais e urbanas.
\end{enumerate}

\section{Limitações}

A despeito dos resultados satisfatórios e da plena validação dos objetivos, foram identificadas limitações inerentes às escolhas de projeto, ao escopo e aos recursos disponíveis:

\begin{enumerate}
    \item \textbf{Limitações Técnicas:} A renderização visual dos elementos no \textit{frontend} utiliza manipulação direta de nós do \textit{Document Object Model} (\gls{dom}) através de \gls{svg} gerenciado pela biblioteca D3.js. Embora ofereça alta interatividade para topologias de porte médio (até algumas centenas de elementos), essa abordagem apresenta degradação de desempenho em termos de taxa de quadros e uso de memória quando submetida a redes gigantescas com milhares de nós simultâneos na mesma tela, situação que exigiria renderização via \textit{canvas} WebGL. Além disso, a aplicação opera em modo local monousuário, sem suporte nativo a múltiplos acessos simultâneos ou banco de dados relacional \textit{multi-tenant}.
    
    \item \textbf{Limitações Metodológicas:} A validação do sistema foi conduzida por meio de simulação lógica e estudos de caso baseados em parâmetros nominais padronizados pelas recomendações ITU-T G.984. O projeto não incluiu a validação experimental de bancada com equipamentos físicos reais (como fontes de luz, medidores de potência óptica ou reflectômetros ópticos no domínio do tempo — \gls{otdr}) operando em bancada de fibra para comparar o sinal simulado com o sinal físico medido.
    
    \item \textbf{Limitações de Escopo:} O escopo da ferramenta concentrou-se exclusivamente na topologia lógica e no orçamento de potência óptica do enlace descendente (\textit{downstream} em 1490 nm), não cobrindo a simulação da dispersão cromática ou do enlace ascendente (\textit{upstream} em 1310 nm). Além disso, o \textit{canvas} visual de modelagem não possui integração nativa com mapas cartográficos georreferenciados (como OpenStreetMap ou Google Maps), exigindo que as distâncias de fibra sejam inseridas manualmente como propriedades dos enlaces. Por fim, o sistema não inclui algoritmos de inteligência computacional ou otimização combinatória para sugestão automática das melhores razões de \textit{splitters} desbalanceados.
    
    \item \textbf{Limitações de Recursos:} O desenvolvimento, testes e medições de desempenho do sistema foram restritos a ambientes computacionais de desenvolvimento pessoal, sem a realização de testes de estresse ou validação em larga escala em servidores de nuvem distribuídos ou em infraestruturas de provedores comerciais em operação.
\end{enumerate}

\section{Trabalhos Futuros}

A partir do conhecimento consolidado e das limitações identificadas neste projeto, vislumbram-se diversas oportunidades de continuidade e aprimoramento do sistema, agrupadas nas seguintes direções:

\begin{enumerate}
    \item \textbf{Extensões Diretas do Sistema:}
    \begin{itemize}
        \item Incorporação de uma camada de mapeamento cartográfico georreferenciado (\gls{gis}) integrada em segundo plano no \textit{canvas} de edição (utilizando bibliotecas como Leaflet ou Mapbox), permitindo calcular automaticamente o comprimento dos cabos a partir das coordenadas geográficas reais de postes e caixas de atendimento.
        \item Evolução da arquitetura \textit{backend} com adição de banco de dados relacional (PostgreSQL com extensão PostGIS), sistema de autenticação de usuários (OAuth2/\gls{jwt}) e suporte ao trabalho colaborativo em tempo real com controle de versões de projetos.
    \end{itemize}
    
    \item \textbf{Novas Direções de Desenvolvimento:}
    \begin{itemize}
        \item Desenvolvimento de um módulo de otimização combinatória inteligente (baseado em algoritmos genéticos ou programação linear inteira) capaz de sugerir automaticamente a combinação ótima das razões dos \textit{splitters} desbalanceados ao longo de um tronco de fibra, visando maximizar a quantidade de atendimentos e garantir que todos os clientes recebam potência dentro da faixa ideal.
        \item Implementação do cálculo óptico bidirecional contemplando a atenuação do canal de \textit{upstream} (1310 nm) e a modelagem do efeito da dispersão cromática e da margem de projeto para sistemas de alta taxa de transmissão.
    \end{itemize}
    
    \item \textbf{Aplicações Práticas e Operacionais:}
    \begin{itemize}
        \item Adoção do \textbf{PON Planner} em fluxos de planejamento técnico e engenharia de redes por Provedores de Serviços de Internet (ISPs) regionais e consultorias especializadas de telecomunicações.
        \item Integração do sistema a ferramentas de monitoramento ativo de redes (como Zabbix ou servidores TR-069 ACS), permitindo comparar continuamente o orçamento óptico simulado na fase de projeto com as medições reais de potência recebida (\textit{Rx Power}) reportadas pelas \glspl{onu} operantes em campo.
    \end{itemize}
    
    \item \textbf{Tecnologias Emergentes:}
    \begin{itemize}
        \item Expansão do motor de cálculo óptico em Python para suportar padrões \gls{pon} de próxima geração (como \gls{xgspon} de 10 Gbps, NG-PON2 e 50G-PON), incorporando as perdas de inserção de filtros de coexistência \gls{wdm} e a análise de múltiplos comprimentos de onda concorrentes (1490 nm, 1577 nm, 1270 nm e 1310 nm).
    \end{itemize}
\end{enumerate}